% !TEX root = ../main.tex
\chapter*{概要}
\addcontentsline{toc}{chapter}{概要}
% ===== 下書き(Claude、11/20)。題目と結論の軸が決まったら直す =====
% 順番:背景 → 課題 → 従来研究の限界 → 提案 → 実装 → 数値つきの結果 → 効かなかった条件

Raspberry Pi などのシングルボードコンピュータを束ねたマイクロクラスタは、安価なエッジ計算基盤として使われている。しかし、専用の空調がない環境に高密度に置かれるため、CPU 温度が上がりやすく、約 80\,℃ を超えるとスロットリングにより性能が落ちる。ファンを常に回せば温度は抑えられるが、ファンの寿命や騒音の点で、必要なときだけ回すことが望ましい。また、囲った環境では、隣のノードの排熱で自ノードの温度が上がる空間熱干渉が生じ、自ノードの負荷だけでは温度を予測できない。

本研究では、隣のノードからの熱影響を状態として推定する拡張熱モデルとカルマンフィルタにより、30 秒先の温度とその不確かさを予測し、不確かさに応じた安全閾値を予測温度が超えるときだけファンを回す制御を提案した。先行発表では、予測に未来の実測値を用いていたことや、同定と評価に同じデータを用いていたことなどの課題があったため、これらを直し、Raspberry Pi 4B 3台と PoE+ HAT のファンを用いて実機で評価した。

その結果、ラックに並べた環境では空間熱干渉はほぼ生じないことが分かった。別のデータで同定したモデルで、負荷が変わらない区間の 30 秒先の温度を平均絶対誤差約 0.7\,℃ で予測できた。ファン制御では、80\,℃ を超えず、スロットリングも起こさずに、常時 ON と比べてファンの稼働を 55〜62\% 減らした。一方、干渉がない環境では予測の不確かさが一定となり、温度だけを見るしきい値方式との差はほぼなかった。PoE+ HAT のファンは効きが速く強いため、予測の利点が小さいと考えられる。
